\chapter{Ingénierie d'Acquisition, Pipeline de Traitement et Intelligence Sémantique}

Dans ce chapitre, nous exposons les mécanismes internes qui transforment la donnée brute, collectée depuis des sources hétérogènes et potentiellement hostiles, en renseignement structuré et actionnable. Nous détaillons successivement l'architecture d'acquisition et ses contraintes de sécurité opérationnelle (OPSEC), le pipeline ETL et sa conformité au standard STIX~2.1, le moteur d'intelligence artificielle, la modélisation avancée par graphe de menace et le modèle de concurrence asynchrone qui garantit la résilience de l'ensemble. Chaque composant est présenté avec la rigueur algorithmique nécessaire à la compréhension de son fonctionnement interne, en cohérence directe avec les choix architecturaux établis au Chapitre~I et les justifications technologiques du Chapitre~II.

%% =========================================================================
\section{Architecture d'Acquisition et Furtivité Opérationnelle (OPSEC)}
%% =========================================================================

L'acquisition constitue la première étape critique de la chaîne de renseignement. Comme établi au Chapitre~I (section~I.4.a), notre matrice d'acquisition opère simultanément sur le Web de surface et sur les couches profondes de l'Internet. Nous présentons ici les mécanismes algorithmiques et les mesures de sécurité opérationnelle (OPSEC) qui régissent chacun de ces vecteurs.

\subsection{Crawlers du Web de Surface : Boucle Événementielle Twisted et Rendu DOM Headless}

Notre pipeline d'ingestion du Web de surface repose sur l'orchestration conjointe de deux cadres complémentaires, conformément à la justification présentée au Chapitre~II (section~II.1.b).

\subsubsection{Scrapy et le paradigme asynchrone Twisted}

Scrapy assure la collecte à grande échelle selon un paradigme asynchrone fondé sur le patron de conception producteur-consommateur. Son architecture interne repose sur la boucle événementielle Twisted (\textit{Twisted Reactor}), une implémentation du patron de conception réacteur (\textit{reactor pattern}) pour Python. Contrairement à un modèle multi-thread dans lequel chaque requête HTTP mobilise un fil d'exécution dédié --- avec le coût mémoire et les risques d'interblocage associés --- Twisted opère sur un fil unique et délègue les opérations d'entrée-sortie réseau au système d'exploitation via des appels non bloquants (\texttt{select}/\texttt{epoll}/\texttt{kqueue} selon la plateforme). L'ordonnanceur central (\texttt{Scheduler}) maintient une file de priorité des requêtes à émettre et les distribue vers un ensemble de téléchargeurs concurrents (\texttt{Downloader}). Chaque réponse HTTP est acheminée vers une fonction d'extraction (\textit{callback}) qui produit soit de nouveaux éléments de données (\texttt{Item}), soit de nouvelles requêtes à planifier (\texttt{Request}). Ce mécanisme garantit une complexité algorithmique sous-linéaire en termes de latence perçue, car les opérations d'entrée-sortie réseau se chevauchent dans le temps sans contention de verrous.

La configuration du moteur de téléchargement applique les contraintes suivantes : un délai inter-requêtes de 2 secondes par domaine (\texttt{DOWNLOAD\_DELAY}), une concurrence maximale de 8 requêtes simultanées par domaine (\texttt{CONCURRENT\_REQUESTS\_PER\_DOMAIN}), et un mécanisme de respect automatique de la directive \texttt{robots.txt} (\texttt{ROBOTSTXT\_OBEY = True}) pour les sources institutionnelles.

\subsubsection{Playwright et le rendu DOM headless}

Playwright intervient comme couche de rendu complémentaire pour les sources qui exigent l'exécution de code JavaScript côté client. Notre implémentation instancie un navigateur Chromium en mode sans interface graphique (\textit{headless}), piloté par le protocole CDP (\textit{Chrome DevTools Protocol}), afin de produire un arbre DOM complet avant l'extraction. Le processus s'exécute en trois phases séquentielles :

\begin{enumerate}
    \item \textbf{Navigation et attente du rendu} : Playwright charge la page cible et attend la stabilisation du réseau (\texttt{networkidle}), ce qui garantit que toutes les requêtes AJAX et les injections dynamiques de contenu sont terminées.
    \item \textbf{Extraction du DOM résolu} : la méthode \texttt{page.content()} retourne le HTML complet tel que rendu par le moteur Blink, incluant les éléments générés dynamiquement par les cadres applicatifs JavaScript (React, Vue, Angular) invisibles dans le HTML source initial.
    \item \textbf{Isolation du processus navigateur} : chaque session de rendu s'exécute dans un contexte de navigateur isolé (\texttt{browser.new\_context()}), empêchant la persistance de cookies, de stockage local ou de données de session entre les collectes.
\end{enumerate}

Cette approche résout le problème des pages à rendu dynamique, fréquentes sur les portails d'actualité en cybersécurité et les tableaux de bord de sources institutionnelles telles que CISA et MITRE ATT\&CK.

\subsubsection{Déduplication par hachage cryptographique SHA-256}

Afin d'éviter la saturation du pipeline par des contenus redondants, nous avons implémenté un intergiciel (\textit{middleware}) de déduplication fondé sur le hachage cryptographique SHA-256. Pour chaque page collectée, l'algorithme calcule une empreinte cryptographique du contenu textuel extrait, après suppression des balises de présentation. Cette empreinte est ensuite comparée à un ensemble de hachages déjà observés, maintenu en mémoire sous forme de table de hachage (\texttt{set} Python) offrant une complexité de recherche en $\mathcal{O}(1)$ amorti.

\begin{lstlisting}[language=Python, caption={Intergiciel de déduplication par hachage SHA-256 avec complexité $\mathcal{O}(1)$.}, label={lst:dedup}]
import hashlib
from bs4 import BeautifulSoup

class ContentDeduplicationMiddleware:
    """
    Intergiciel Scrapy : rejette les pages dont le contenu
    textuel a deja ete observe. Complexite O(1) amorti
    via table de hachage (set Python).
    """
    def __init__(self):
        self._seen_hashes: set[str] = set()

    def process_response(self, request, response, spider):
        soup = BeautifulSoup(response.text, "html.parser")
        # Suppression des elements non informatifs
        for tag in soup(["script", "style", "noscript",
                         "nav", "footer", "header"]):
            tag.decompose()
        clean_text = soup.get_text(separator="\n", strip=True)

        # Calcul de l'empreinte SHA-256
        content_hash = hashlib.sha256(
            clean_text.encode("utf-8", errors="ignore")
        ).hexdigest()

        # Test d'appartenance en O(1) amorti
        if content_hash in self._seen_hashes:
            spider.logger.debug(
                f"Doublon detecte, rejet : {request.url}"
            )
            from scrapy.exceptions import IgnoreRequest
            raise IgnoreRequest("Contenu duplique")

        self._seen_hashes.add(content_hash)
        response.meta["content_hash"] = content_hash
        return response
\end{lstlisting}

Formellement, pour un document $d_i$ dont le contenu textuel nettoyé est $T_i$, nous calculons $h_i = \texttt{SHA-256}(T_i)$. Si $h_i \in \mathcal{H}_{\text{vu}}$, le document est rejeté comme doublon. Sinon, $h_i$ est inséré dans $\mathcal{H}_{\text{vu}}$ et le document progresse vers l'étape suivante du pipeline. La propriété de résistance aux collisions de SHA-256 --- $2^{128}$ opérations pour une collision par attaque par anniversaire --- assure que deux contenus distincts ne produisent jamais la même empreinte dans le cadre opérationnel de notre plateforme.

\subsection{Infiltration du Dark Web : Gestion du Routage en Oignon et OPSEC}

La collecte sur les couches profondes de l'Internet constitue le vecteur d'acquisition le plus sensible de notre architecture. Nous avons conçu un gestionnaire de circuits Tor (\texttt{TorManager}) qui encapsule l'intégralité du cycle de vie des connexions au réseau d'anonymisation, depuis l'établissement initial du circuit jusqu'à la rotation automatisée des adresses IP de sortie.

\subsubsection{Automate d'état du gestionnaire Tor}

Notre gestionnaire opère selon un automate à quatre états : \texttt{CONNECTED}, \texttt{DISCONNECTED}, \texttt{ROTATING} et \texttt{KILLED}. L'état \texttt{KILLED} correspond à l'activation d'un mécanisme d'arrêt d'urgence (\textit{kill-switch}) qui bloque instantanément tout trafic vers les services cachés \texttt{.onion}. Cette conception défensive garantit qu'en cas de compromission suspectée du circuit, l'opérateur peut immédiatement couper toute interaction avec les sources profondes, sans attendre la terminaison des connexions en cours.

La connectivité est vérifiée par interrogation de l'interface programmatique du Tor Project (\texttt{check.torproject.org/api/ip}), qui confirme que le trafic est effectivement acheminé à travers le réseau Tor et retourne l'adresse IP du n\oe ud de sortie courant. Cette vérification est systématiquement exécutée avant chaque session de collecte.

\subsubsection{Rotation des circuits via le signal NEWNYM (port de contrôle 9051)}

Nous avons implémenté la rotation automatisée des circuits Tor par communication directe avec le processus Tor via son port de contrôle (port~9051). Le signal \texttt{NEWNYM} ordonne au client Tor de construire de nouveaux circuits pour les requêtes futures, modifiant ainsi l'adresse IP de sortie observable par les sites cibles. Notre implémentation utilise le protocole de contrôle Tor sur une connexion TCP authentifiée, avec deux critères de déclenchement complémentaires : un seuil de 50 requêtes par circuit et un intervalle temporel maximal configurable.

\begin{lstlisting}[language=Python, caption={Rotation de circuit Tor via signal NEWNYM sur le port de contrôle 9051.}, label={lst:newnym}]
import asyncio
import random
import math

class TorCircuitRotator:
    """
    Gestionnaire de rotation de circuits Tor.
    Envoie le signal NEWNYM au port de contrôle 9051
    apres verification de la condition de seuil.
    """
    def __init__(self, control_port=9051,
                 control_password="onyx_tor_auth",
                 max_requests_per_circuit=50):
        self._control_port = control_port
        self._password = control_password
        self._max_requests = max_requests_per_circuit
        self._request_count = 0
        self._lock = asyncio.Lock()

    def _stochastic_delay(self) -> float:
        """
        Delai stochastique par distribution log-normale
        tronquee : mu=1.5s, sigma=0.5s, borne [0.8, 6.0].
        Emule un comportement humain non deterministe.
        """
        mu, sigma = 1.5, 0.5
        delay = random.lognormvariate(
            math.log(mu), sigma
        )
        return max(0.8, min(delay, 6.0))

    async def maybe_rotate(self):
        """Rotation conditionnelle apres N requetes."""
        self._request_count += 1
        # Condition de declenchement
        if self._request_count < self._max_requests:
            return
        await self._send_newnym()

    async def _send_newnym(self):
        """Envoi du signal NEWNYM au processus Tor."""
        async with self._lock:
            reader, writer = await asyncio.wait_for(
                asyncio.open_connection(
                    "127.0.0.1", self._control_port
                ),
                timeout=10.0,
            )
            await reader.readline()  # Banniere Tor

            # Authentification au port de controle
            writer.write(
                f'AUTHENTICATE "{self._password}"\r\n'
                .encode()
            )
            await writer.drain()
            resp = await reader.readline()
            if b"250" not in resp:
                raise ConnectionError(
                    f"Echec authentification Tor: "
                    f"{resp.decode().strip()}"
                )

            # Emission du signal NEWNYM
            writer.write(b"SIGNAL NEWNYM\r\n")
            await writer.drain()
            resp = await reader.readline()

            writer.close()
            await writer.wait_closed()

            if b"250" in resp:
                # Attente de construction du circuit
                delay = self._stochastic_delay()
                await asyncio.sleep(delay)
                self._request_count = 0
\end{lstlisting}

Le délai post-rotation est calculé par une distribution log-normale tronquée $\mathcal{LN}(\mu=1{,}5, \sigma=0{,}5)$ bornée à l'intervalle $[0{,}8; 6{,}0]$ secondes. Nous avons retenu la distribution log-normale plutôt qu'une distribution uniforme, car elle modélise plus fidèlement les temps de réaction humains : une densité de probabilité concentrée autour de la médiane avec une queue droite étendue. Ce profil temporel non déterministe empêche l'identification du comportement du collecteur par analyse statistique des temps inter-requêtes, un vecteur de détection classique des systèmes anti-automatisation.

\subsubsection{Contournement de l'anti-fingerprinting des blogs de rançongiciels}

Les blogs de rançongiciels et les forums clandestins déploient des mécanismes de détection des agents automatisés fondés sur l'analyse de l'empreinte du navigateur (\textit{browser fingerprint}). Notre architecture contourne ces dispositifs par trois mesures complémentaires de ségrégation :

\begin{enumerate}
    \item \textbf{Rotation des agents utilisateur} : chaque requête est émise avec un en-tête \texttt{User-Agent} sélectionné aléatoirement parmi un ensemble de cinq signatures correspondant aux versions authentiques du navigateur Tor Browser (basé sur Firefox~ESR~115). Le trafic devient ainsi indistinguable du trafic légitime des utilisateurs du réseau Tor.
    \item \textbf{Conformité protocolaire stricte} : les en-têtes HTTP (\texttt{Accept}, \texttt{Accept-Language}, \texttt{Accept-Encoding}, \texttt{Connection}, \texttt{Upgrade-Insecure-Requests}) reproduisent exactement le profil par défaut du navigateur Tor Browser, évitant toute divergence détectable par les systèmes d'analyse d'en-têtes côté serveur.
    \item \textbf{Isolation par domaine (IsolateDestAddr)} : la directive \texttt{IsolateDestAddr} dans la configuration \texttt{torrc} garantit que chaque domaine \texttt{.onion} cible est atteint via un circuit Tor distinct. Ce cloisonnement empêche la corrélation croisée entre les sessions de collecte sur des cibles différentes, même si un n\oe ud de sortie compromis observe le trafic.
\end{enumerate}

La conjonction de ces mesures assure une furtivité opérationnelle suffisante pour opérer de manière soutenue sur les principales places de marché illicites et les blogs de publication de données exfiltrées, sans déclencher les mécanismes de bannissement par adresse IP ou par empreinte comportementale.

%% =========================================================================
\section{Pipeline ETL (Extract, Transform, Load) et Modélisation STIX~2.1}
%% =========================================================================

La donnée brute extraite par la matrice d'acquisition se présente sous une forme hétérogène : pages HTML, flux JSON, contenus textuels non structurés issus de forums clandestins. Nous avons conçu un pipeline ETL (\textit{Extract, Transform, Load}) rigoureux qui normalise, transforme et persiste cette donnée selon le standard STIX~2.1, garantissant ainsi l'interopérabilité exigée par l'exigence fonctionnelle EF-2 du Chapitre~I.

\subsection{Phase d'extraction : Nettoyage et Refangage}

La phase d'extraction opère en deux temps. Le parseur HTML, implémenté via BeautifulSoup4 (Chapitre~II, section~II.1.b), applique d'abord une décomposition sélective de l'arbre DOM : les éléments non informatifs (\texttt{<script>}, \texttt{<style>}, \texttt{<noscript>}, \texttt{<nav>}, \texttt{<footer>}, \texttt{<header>}) sont systématiquement supprimés avant l'extraction du texte résiduel. Cette opération réduit en moyenne de 60~\% à 80~\% le volume de données transmis aux étapes analytiques en aval, diminuant proportionnellement la charge computationnelle des algorithmes de traitement du langage naturel.

Le texte nettoyé est ensuite soumis à un processus de refangage (\textit{refanging}) par application séquentielle de sept expressions régulières précompilées. Les indicateurs de compromission intentionnellement obscurcis par leurs auteurs --- convention courante dans les rapports techniques et les forums de cybersécurité --- sont automatiquement restaurés dans leur forme canonique :

\begin{itemize}
    \item \texttt{hxxp} $\rightarrow$ \texttt{http} et \texttt{hXXp} $\rightarrow$ \texttt{http}
    \item \texttt{[.]} et \texttt{(.)} et \texttt{[dot]} $\rightarrow$ \texttt{.}
    \item \texttt{[://]} $\rightarrow$ \texttt{://}
    \item \texttt{[at]} et \texttt{(at)} $\rightarrow$ \texttt{@}
    \item \texttt{\textbackslash.} $\rightarrow$ \texttt{.}
\end{itemize}

Ce processus est systématiquement appliqué en amont de la phase d'extraction d'entités, garantissant que les algorithmes de reconnaissance opèrent sur des représentations normalisées indépendamment des conventions de défangage propres à chaque source.

\subsection{Phase de transformation : Mapping vers le standard STIX~2.1}

La transformation constitue l'étape pivot du pipeline ETL. Chaque indicateur de compromission extrait est mappé vers un objet STIX~2.1 de type \texttt{indicator}, dont le champ \texttt{pattern} est construit selon la grammaire formelle STIX Patterning définie par le comité technique OASIS CTI. Notre implémentation supporte huit types d'indicateurs avec leur patron correspondant : adresses IPv4 et IPv6, noms de domaine, URL, empreintes cryptographiques (MD5, SHA-1, SHA-256) et adresses de courrier électronique.

L'objet STIX généré par notre pipeline se présente sous la forme JSON suivante :

\begin{lstlisting}[language=json, caption={Objet STIX~2.1 de type \texttt{indicator} généré par le pipeline ONYX.}, label={lst:stix-json}]
{
  "type": "indicator",
  "spec_version": "2.1",
  "id": "indicator--a1b2c3d4-e5f6-7890-abcd-ef1234567890",
  "created": "2026-05-12T14:30:00.000Z",
  "modified": "2026-05-12T14:30:00.000Z",
  "name": "ipv4: 185.220.101.45",
  "description": "C2 server detected by ONYX NLP pipeline",
  "pattern": "[ipv4-addr:value = '185.220.101.45']",
  "pattern_type": "stix",
  "valid_from": "2026-05-12T14:30:00.000Z",
  "confidence": 97,
  "labels": ["ipv4", "abuse.ch Feodo Tracker"],
  "x_onyx_source": "abuse.ch Feodo Tracker",
  "x_onyx_tags": ["feodo", "c2", "botnet"],
  "x_onyx_mitre_techniques": ["T1071", "T1568"],
  "x_onyx_decay_score": 1.0,
  "x_onyx_decay_state": "valid"
}
\end{lstlisting}

Nous soulignons plusieurs choix de conception dans cette structure. Le champ \texttt{pattern} respecte strictement la grammaire STIX Patterning, ce qui garantit l'interopérabilité avec les plateformes tierces (MISP, OpenCTI, Splunk). Les champs préfixés par \texttt{x\_onyx\_} exploitent le mécanisme d'extension du standard STIX~2.1 pour enrichir chaque indicateur avec des métadonnées propriétaires --- source d'acquisition, étiquettes de classification, techniques MITRE associées, score de déclin --- sans violer la conformité au standard. Le champ \texttt{confidence} (0--100) quantifie la fiabilité de l'indicateur selon un barème combinant la réputation de la source et le score de confiance de l'extracteur NLP.

\subsection{Phase de chargement : Persistance MongoDB et le problème du polymorphisme STIX}

\subsubsection{Le problème du polymorphisme STIX dans un schéma relationnel}

Le modèle de données STIX~2.1 définit 18 types d'objets de domaine (SDO), 2 types d'objets de relation (SRO) et 6 types de méta-objets, chacun possédant un ensemble d'attributs qui lui est propre. Un objet \texttt{indicator} possède les champs \texttt{pattern}, \texttt{pattern\_type} et \texttt{valid\_from}; un objet \texttt{threat-actor} possède les champs \texttt{threat\_actor\_types}, \texttt{aliases} et \texttt{goals}; un objet \texttt{malware} possède les champs \texttt{malware\_types}, \texttt{is\_family} et \texttt{capabilities}. Ce polymorphisme structurel pose un problème fondamental pour la persistance relationnelle.

Dans une base relationnelle telle que PostgreSQL, trois stratégies sont envisageables pour modéliser ce polymorphisme, et chacune présente des limitations critiques :

\begin{itemize}
    \item \textbf{Table unique avec colonnes nullables} (\textit{Single Table Inheritance}) : une table monolithique contient l'union de tous les attributs possibles. Pour 18 types d'objets, cela produit une table excessivement sparse avec des dizaines de colonnes à valeur nulle pour chaque enregistrement, dégradant l'efficacité de stockage et la lisibilité du schéma.
    \item \textbf{Tables par type} (\textit{Concrete Table Inheritance}) : chaque type STIX possède sa propre table. Les requêtes transversales (par exemple, ``tous les objets créés depuis 24 heures'') nécessitent des opérations \texttt{UNION} sur 18 tables, dont la complexité croît linéairement avec le nombre de types et le volume de chaque table.
    \item \textbf{Modèle Entité-Attribut-Valeur (EAV)} : les attributs sont stockés sous forme de triplets $(entité, attribut, valeur)$. Bien que flexible, ce modèle transforme chaque lecture d'objet en une série de jointures dont le nombre est proportionnel au nombre d'attributs, produisant une charge I/O-bound prohibitive à grande échelle.
\end{itemize}

\subsubsection{MongoDB et la persistance documentaire naturelle}

MongoDB, en tant que base documentaire, résout ce problème de polymorphisme par construction. Chaque objet STIX est stocké comme un document JSON autonome dont la structure s'adapte naturellement à son type, sans contrainte de schéma fixe. Un document \texttt{indicator} contient exactement les champs de l'indicateur; un document \texttt{threat-actor} contient exactement les champs de l'acteur de menace. Aucune colonne nulle, aucune jointure, aucune opération \texttt{UNION}.

Notre service MongoDB implémente le patron Singleton avec un pool de connexions asynchrones via Motor (wrapper asynchrone de PyMongo), configuré avec un minimum de 5 et un maximum de 50 connexions simultanées (\texttt{minPoolSize=5}, \texttt{maxPoolSize=50}). Les temporisations sont calibrées pour l'environnement conteneurisé : 5~secondes pour la sélection du serveur, 5~secondes pour l'établissement de la connexion, et 30~secondes pour les opérations de lecture sur socket.

\subsubsection{Insertion idempotente par upsert}

Les opérations d'écriture utilisent systématiquement le mécanisme d'insertion idempotente (\textit{upsert}) de MongoDB. Pour chaque objet STIX reçu, la méthode \texttt{update\_one()} est invoquée avec le filtre \texttt{\{"id": stix\_obj.id\}} et l'option \texttt{upsert=True}. Si un document portant le même identifiant STIX existe déjà, il est mis à jour avec les nouvelles métadonnées (notamment l'horodatage \texttt{modified}). Sinon, un nouveau document est créé. Ce mécanisme garantit qu'un même indicateur reçu simultanément de sources multiples --- situation fréquente lorsque plusieurs flux OSINT convergent vers le même IoC --- ne produit pas de doublons dans la base, tout en assurant que chaque observation successive enrichit le document existant.

Pour les opérations de masse, nous avons implémenté une méthode \texttt{bulk\_create\_stix()} qui regroupe les documents par collection cible (objets STIX, relations, observations) et exécute un \texttt{bulk\_write()} avec ordonnancement désactivé (\texttt{ordered=False}), ce qui permet au pilote de paralléliser les écritures et de poursuivre l'insertion en cas d'erreur sur un document individuel, sans interrompre le lot complet.

Le routage vers la collection appropriée est assuré par une fonction déterministe qui inspecte le champ \texttt{type} de chaque objet : les objets de type \texttt{relationship} sont dirigés vers la collection \texttt{stix\_relationships}, les objets de type \texttt{sighting} vers \texttt{stix\_sightings}, les définitions de marquage vers \texttt{marking\_definitions}, et tous les autres types vers la collection générale \texttt{stix\_objects}. Cette ségrégation par collection optimise les performances de requêtage en réduisant l'espace de recherche pour les opérations les plus fréquentes : la traversée de relations et la recherche d'indicateurs.



\section{Le C\oe ur Cognitif : Implémentation de l'Intelligence Artificielle}

Le moteur cognitif constitue le noyau analytique de la plateforme, tel que défini au Chapitre~I (section~I.4.b). Sa mission est double : extraire automatiquement les entités de renseignement enfouies dans des flux textuels hétérogènes, puis projeter ces informations dans un espace vectoriel de haute dimension où la proximité géométrique encode la similarité sémantique. Nous détaillons ici les deux fonctions algorithmiques centrales qui réalisent cette transformation, en exposant leurs fondements mathématiques et leurs garanties de performance.

\subsection{Extraction d'Entités Nommées (NER) et Validations Strictes}

Notre moteur d'extraction d'indicateurs de compromission (\texttt{IOCExtractor}) implémente une architecture en deux phases --- extraction large puis validation stricte --- inspirée des patrons de conception du cadriciel AIL (développé par le CERT luxembourgeois) et de la plateforme IntelOwl. Cette conception garantit un taux de faux positifs proche de zéro tout en maintenant un rappel élevé sur l'ensemble des 18 types d'indicateurs supportés.

\subsubsection{Cascade d'expressions régulières par spécificité décroissante}

La première phase déploie un ensemble de 18 expressions régulières précompilées (\texttt{re.compile}), chacune ciblant un type spécifique d'indicateur. L'ordre d'application de ces patrons n'est pas arbitraire : il obéit à un principe de \textbf{précédence par spécificité décroissante} dont la violation produirait des erreurs de classification systématiques.

Considérons le cas des empreintes cryptographiques. Un hachage SHA-512 est une chaîne de 128 caractères hexadécimaux. Un hachage SHA-256 en comporte 64, un SHA-1 en comporte 40 et un MD5 en comporte 32. Si le patron MD5 (32 caractères hexadécimaux contigus) était appliqué en premier, les 32 premiers caractères d'un hachage SHA-256 seraient capturés comme un faux MD5, et le hachage SHA-256 légitime ne serait jamais identifié. Formellement, pour deux patrons $P_a$ et $P_b$ où $\mathcal{L}(P_a) \subset \mathcal{L}(P_b)$ (le langage reconnu par $P_a$ est un sous-ensemble strict du langage reconnu par $P_b$), le patron $P_a$ --- le plus restrictif --- doit être appliqué en premier.

Notre implémentation applique donc la séquence suivante :

\begin{enumerate}
    \item SHA-512 (128 caractères hexadécimaux) --- le plus restrictif.
    \item SHA-256 (64 caractères) --- tout SHA-256 résiduel après élimination des SHA-512.
    \item SHA-1 (40 caractères) --- extraction après élimination des correspondances plus longues.
    \item MD5 (32 caractères) --- le moins restrictif des hachages.
    \item Empreintes JA3/JA3S (32 caractères avec préfixe \texttt{ja3:} ou \texttt{ja3s:}).
    \item Identifiants CVE, techniques MITRE ATT\&CK.
    \item URL (avant les domaines, car les URL contiennent des domaines).
    \item Adresses de courrier électronique (avant les domaines, pour la même raison).
    \item Plages CIDR (avant les adresses IP individuelles).
    \item Adresses IPv4 et IPv6.
    \item Noms de domaine.
    \item Portefeuilles de cryptomonnaies (Bitcoin, Ethereum, Monero).
    \item Adresses de services cachés \texttt{.onion}, numéros ASN, règles YARA.
\end{enumerate}

Un mécanisme de déduplication inter-patrons complète cette cascade : avant d'enregistrer une correspondance, l'algorithme vérifie que la valeur extraite n'est pas déjà contenue dans une correspondance plus longue précédemment capturée. Cette vérification empêche, par exemple, qu'un sous-ensemble de 32 caractères d'un SHA-256 déjà extrait soit enregistré comme un MD5 distinct.

\subsubsection{Filtrage topologique : exclusion des plages réservées RFC~1918}

Chaque adresse IPv4 candidate est soumise à un filtrage topologique rigoureux avant son admission comme indicateur de compromission. Ce filtrage repose sur le module \texttt{ipaddress} de la bibliothèque standard Python, qui fournit une implémentation conforme aux RFC pour la manipulation d'adresses réseau.

Nous excluons systématiquement les adresses appartenant aux plages suivantes :

\begin{itemize}
    \item \textbf{Plages privées RFC~1918} : \texttt{10.0.0.0/8}, \texttt{172.16.0.0/12}, \texttt{192.168.0.0/16} --- ces adresses ne sont pas routables sur l'Internet public et ne constituent donc jamais des indicateurs de compromission valides dans un contexte de renseignement sur les menaces externes.
    \item \textbf{Adresses de bouclage} : \texttt{127.0.0.0/8} --- adresses locales sans signification opérationnelle.
    \item \textbf{Adresses de lien local} : \texttt{169.254.0.0/16} --- adresses auto-attribuées en l'absence de serveur DHCP.
    \item \textbf{Plages de multidiffusion et réservées} : \texttt{224.0.0.0/4} (multidiffusion) et \texttt{240.0.0.0/4} (réservé pour usage futur).
    \item \textbf{Plage nulle} : \texttt{0.0.0.0/8} --- adresses non attribuables.
\end{itemize}

Le test d'appartenance est effectué par itération sur un ensemble préchargé de huit objets \texttt{ip\_network}, avec une complexité de $\mathcal{O}(1)$ par test (nombre fixe de réseaux) et une complexité totale de $\mathcal{O}(k)$ par adresse candidate, où $k = 8$ est le nombre de plages exclues.

\subsubsection{Analyse entropique et rejet des faux positifs}

Les empreintes cryptographiques légitimes présentent une entropie informationnelle élevée : la sortie d'une fonction de hachage cryptographique est, par construction, pseudo-aléatoire et uniformément distribuée sur l'espace de sortie. À l'inverse, les chaînes hexadécimales accidentellement capturées par nos expressions régulières --- identifiants de version, clés de configuration, fragments de code --- présentent typiquement une entropie significativement plus faible.

Nous exploitons cette propriété en calculant la cardinalité de l'alphabet effectif de chaque correspondance candidate. Pour une chaîne $s$ de longueur $\ell$ composée de caractères hexadécimaux (alphabet théorique $\Sigma = \{0, 1, \ldots, 9, a, b, c, d, e, f\}$, $|\Sigma| = 16$), nous définissons la diversité caractérielle $\delta(s) = |\{c \in \Sigma : c \in s\}|$. Un seuil minimal $\delta_{\min} = 6$ est imposé : toute chaîne dont la diversité est inférieure à 6 caractères hexadécimaux distincts est rejetée comme faux positif.

Ce seuil a été calibré empiriquement : une empreinte SHA-256 authentique, produite par une fonction pseudo-aléatoire, contient en moyenne 15{,}4 caractères distincts sur 16 possibles pour une chaîne de 64 caractères (d'après le problème du collecteur de coupons, $\mathbb{E}[\delta] = 16 \cdot (1 - (15/16)^{64}) \approx 15{,}4$). Un seuil de 6 offre donc une marge confortable de rejet des faux positifs tout en éliminant les chaînes pathologiques telles que \texttt{0000...0000} ou \texttt{aaaa...aaaa}.

\begin{lstlisting}[language=Python, caption={Validation d'empreinte cryptographique par analyse entropique et exclusion contextuelle.}, label={lst:entropy}]
import re

# Empreintes triviales connues (faux positifs certains)
HASH_EXCLUSIONS = {
    "0" * 32, "f" * 32,   # MD5 nuls ou pleins
    "0" * 40, "f" * 40,   # SHA-1 nuls ou pleins
    "0" * 64, "f" * 64,   # SHA-256 nuls ou pleins
    "0" * 128, "f" * 128, # SHA-512 nuls ou pleins
}

# Contexte indicatif d'un faux positif
FP_CONTEXT_PATTERN = re.compile(
    r"version|v\d|release|update|build",
    re.IGNORECASE,
)

def is_valid_hash(
    hash_str: str,
    expected_len: int,
    context: str = "",
    entropy_threshold: int = 6,
) -> bool:
    """
    Validation multi-criteres d'une empreinte
    cryptographique candidate.

    Criteres :
      1. Longueur exacte attendue.
      2. Non-appartenance aux empreintes triviales.
      3. Diversite characterielle >= entropy_threshold.
      4. Absence de contexte de faux positif.

    Retourne True si l'empreinte est consideree valide.
    """
    # Critere 1 : longueur stricte
    if len(hash_str) != expected_len:
        return False

    # Critere 2 : exclusion des empreintes triviales
    if hash_str.lower() in HASH_EXCLUSIONS:
        return False

    # Critere 3 : analyse entropique
    # delta(s) = |{c in Sigma : c in s}|
    unique_chars = len(set(hash_str.lower()))
    if unique_chars < entropy_threshold:
        return False

    # Critere 4 : analyse contextuelle
    # Rejet si le contexte environnant (+/- 50 chars)
    # contient des indicateurs de version/release
    if context and FP_CONTEXT_PATTERN.search(context):
        return False

    return True
\end{lstlisting}

\subsubsection{Analyse contextuelle et scoring de confiance}

Au-delà du filtrage binaire (acceptation ou rejet), notre extracteur attribue à chaque indicateur un score de confiance $\gamma \in [0, 1]$ qui quantifie la probabilité que la correspondance soit un véritable indicateur de compromission. Ce score est initialisé à une valeur par défaut dépendant du type (0{,}95 pour les adresses IPv4 publiques, 0{,}90 pour les empreintes cryptographiques, 0{,}85 pour les noms de domaine) puis modulé par l'analyse contextuelle.

Pour chaque correspondance, une fenêtre de contexte de $\pm 50$ caractères est extraite du texte source. Si cette fenêtre contient des termes indicatifs d'un contexte non malveillant --- \textit{version}, \textit{release}, \textit{update}, \textit{build} --- le score de confiance est abaissé de 0{,}95 à 0{,}40. Les correspondances dont le score final est inférieur au seuil minimal configurable ($\gamma_{\min} = 0{,}50$ par défaut) sont exclues du résultat final. Cette approche graduée offre une flexibilité opérationnelle supérieure au filtrage binaire : l'analyste peut ajuster le seuil de sensibilité en fonction du contexte opérationnel, privilégiant le rappel (seuil bas) ou la précision (seuil haut) selon les besoins de sa mission.

\subsection{Vectorisation Sémantique (SciBERT) et Indexation HNSW (Qdrant)}

La recherche par mots-clés, même enrichie par les opérateurs booléens et la pondération TF-IDF d'Elasticsearch (Chapitre~II, section~II.3.d), ne suffit pas à détecter les similarités conceptuelles entre rapports de renseignement. Deux rapports décrivant la même campagne d'attaque peuvent utiliser des vocabulaires radicalement différents : l'un parlera d'\textit{exfiltration via DNS tunneling}, l'autre de \textit{canal de commande et contrôle exploitant des requêtes DNS}. La correspondance lexicale entre ces deux formulations est nulle, alors que leur proximité sémantique est maximale. Seule une représentation dans un espace vectoriel de haute dimension permet de capturer cette équivalence conceptuelle.

\subsubsection{Architecture de SciBERT et spécialisation sur corpus scientifique}

SciBERT est un modèle de langage pré-entraîné de la famille BERT (\textit{Bidirectional Encoder Representations from Transformers}), développé par l'Allen Institute for AI. Contrairement au modèle BERT généraliste --- entraîné sur les corpus BooksCorpus et Wikipedia anglophone --- SciBERT a été pré-entraîné sur un corpus de 1{,}14 million d'articles scientifiques issus de Semantic Scholar, couvrant les domaines de l'informatique et des sciences biomédicales. Ce pré-entraînement spécialisé confère au modèle une compréhension fine du vocabulaire technique utilisé dans les rapports de cybersécurité : termes tels que \textit{lateral movement}, \textit{command and control}, \textit{privilege escalation} ou \textit{ransomware-as-a-service} sont correctement tokenisés et contextualisés, là où un modèle généraliste les fragmenterait en sous-mots peu informatifs.

L'architecture interne de SciBERT se compose de :
\begin{itemize}
    \item \textbf{Couche de segmentation (WordPiece)} : le texte d'entrée est décomposé en sous-mots (\textit{subword tokens}) via l'algorithme WordPiece, qui maintient un vocabulaire de 31~000 sous-mots appris sur le corpus scientifique. Les termes techniques fréquents (\textit{malware}, \textit{vulnerability}) sont conservés comme jetons uniques; les termes rares sont décomposés en fragments connus du vocabulaire.
    \item \textbf{12 couches de Transformeur} : chaque couche applique un mécanisme d'attention multi-têtes (\textit{multi-head self-attention}) composé de 12 têtes d'attention indépendantes. Pour une séquence d'entrée $\mathbf{X} \in \mathbb{R}^{n \times 768}$ (où $n$ est le nombre de jetons), chaque tête $h$ calcule :
    \[
    \text{Attention}_h(\mathbf{Q}_h, \mathbf{K}_h, \mathbf{V}_h) = \text{softmax}\!\left(\frac{\mathbf{Q}_h \mathbf{K}_h^\top}{\sqrt{d_k}}\right) \mathbf{V}_h
    \]
    où $\mathbf{Q}_h = \mathbf{X}\mathbf{W}_h^Q$, $\mathbf{K}_h = \mathbf{X}\mathbf{W}_h^K$, $\mathbf{V}_h = \mathbf{X}\mathbf{W}_h^V$ sont les matrices de requête, clé et valeur, et $d_k = 64$ est la dimension de chaque tête. Les sorties des 12 têtes sont concaténées et projetées linéairement pour produire la sortie de la couche. Ce mécanisme permet au modèle de pondérer l'importance relative de chaque jeton en fonction de l'ensemble du contexte, capturant les dépendances sémantiques à longue distance.
    \item \textbf{Couche de sortie} : le vecteur associé au jeton spécial \texttt{[CLS]}, inséré en position initiale de chaque séquence, encode une représentation agrégée de l'ensemble du document après traversée des 12 couches d'attention. Ce vecteur $\mathbf{v}_r \in \mathbb{R}^{768}$ constitue la représentation sémantique du rapport $r$.
\end{itemize}

\subsubsection{Similarité cosinus et corrélation sémantique}

La mesure de similarité retenue pour comparer les représentations vectorielles est la similarité cosinus. Pour deux vecteurs $\mathbf{u}, \mathbf{v} \in \mathbb{R}^{768}$ représentant respectivement deux rapports de renseignement, elle est définie par :

\begin{equation}
\label{eq:cosine-sim}
\text{sim}_{\cos}(\mathbf{u}, \mathbf{v}) = \frac{\mathbf{u} \cdot \mathbf{v}}{\|\mathbf{u}\|_2 \; \|\mathbf{v}\|_2} = \frac{\displaystyle\sum_{i=1}^{768} u_i \, v_i}{\sqrt{\displaystyle\sum_{i=1}^{768} u_i^2} \;\; \sqrt{\displaystyle\sum_{i=1}^{768} v_i^2}}
\end{equation}

Cette mesure possède deux propriétés mathématiques qui la rendent supérieure à la recherche lexicale pour la corrélation de campagnes d'attaques :

\begin{enumerate}
    \item \textbf{Invariance à la norme} : la similarité cosinus mesure exclusivement l'angle entre les directions des vecteurs dans $\mathbb{R}^{768}$, indépendamment de leur magnitude. Deux rapports de longueurs très différentes --- un bulletin d'alerte de 200 mots et un rapport d'investigation de 5\,000 mots --- décrivant la même campagne produiront des vecteurs de normes différentes mais orientés dans la même direction, et obtiendront donc une similarité élevée. La recherche par mots-clés (TF-IDF, BM25) est au contraire biaisée par la fréquence absolue des termes, pénalisant les documents courts.
    \item \textbf{Capture des relations sémantiques latentes} : les couches d'attention de SciBERT projettent les synonymes, les paraphrases et les formulations alternatives vers des régions proches de l'espace vectoriel. Les termes \textit{data exfiltration} et \textit{information theft}, bien que lexicalement disjoints, produisent des représentations vectorielles dont la similarité cosinus est typiquement supérieure à 0{,}90 après pré-entraînement. La recherche Elasticsearch, fondée sur la correspondance de jetons (\textit{token matching}) et la distance d'édition (\textit{fuzzy matching}), ne peut pas capturer cette équivalence sémantique car elle opère dans l'espace des formes lexicales de surface, et non dans l'espace des significations.
\end{enumerate}

En pratique, nous définissons un seuil de corrélation $\tau = 0{,}85$ : toute paire de rapports dont la similarité cosinus dépasse ce seuil est signalée comme potentiellement liée à la même campagne, même en l'absence totale de mots-clés communs. Cette capacité de corrélation sémantique constitue un avantage fonctionnel déterminant par rapport aux solutions analysées au Chapitre~I (section~I.2) : ni MISP, ni OpenCTI ne proposent nativement de recherche par similarité vectorielle.

\subsubsection{Indexation HNSW dans Qdrant et complexité logarithmique}

Les vecteurs produits par SciBERT sont indexés dans Qdrant, notre base de données vectorielle (Chapitre~II, section~II.3.e). Le défi algorithmique de la recherche du plus proche voisin (\textit{Nearest Neighbor Search}) dans un espace de dimension $d = 768$ est le suivant : une recherche exhaustive (\textit{brute-force}) nécessite le calcul de la similarité cosinus entre le vecteur requête et chacun des $n$ vecteurs indexés, soit une complexité de $\mathcal{O}(n \cdot d)$. Pour un corpus de $n = 10^6$ rapports, cette approche est prohibitive en temps de réponse interactif.

Qdrant résout ce problème par l'indexation HNSW (\textit{Hierarchical Navigable Small World}), une structure de données probabiliste qui organise les vecteurs en un graphe navigable multi-couches. L'algorithme HNSW construit une hiérarchie de graphes de proximité : la couche supérieure contient un sous-ensemble réduit de n\oe uds avec des arêtes à longue portée pour la navigation grossière, tandis que les couches inférieures ajoutent progressivement des n\oe uds avec des arêtes à courte portée pour le raffinement local. La recherche s'effectue par descente gloutonne (\textit{greedy search}) depuis la couche supérieure vers la couche de base.

La complexité de recherche résultante est en $\mathcal{O}(\log n)$ pour $n$ vecteurs indexés, soit un gain de plusieurs ordres de grandeur par rapport à la recherche exhaustive. Pour un corpus de $10^6$ documents, le nombre d'opérations de comparaison passe de $10^6$ (recherche exhaustive) à environ 20 (recherche HNSW), avec une dégradation contrôlée de l'exactitude : le paramètre \texttt{ef\_search} permet de régler le compromis entre précision et latence, un \texttt{ef\_search} de 128 offrant typiquement un rappel supérieur à 99~\% des vrais plus proches voisins.

Cette architecture vectorielle complète la recherche textuelle d'Elasticsearch selon un modèle de complémentarité fonctionnelle : Elasticsearch excelle dans la recherche exacte de termes connus (\textit{``CVE-2024-3400''}, \textit{``185.220.101.45''}), tandis que Qdrant excelle dans la découverte de relations sémantiques latentes entre documents dont le vocabulaire de surface diverge. Nous exploitons conjointement ces deux capacités dans le pipeline RAG (\textit{Retrieval-Augmented Generation}) qui alimente l'assistant conversationnel de la plateforme.



\section{Modélisation Avancée : Graphe de Menace et Dégradation Tactique}

Au-delà de l'analyse individuelle des indicateurs, la valeur opérationnelle du renseignement réside dans la capacité à reconstruire les structures relationnelles qui lient les entités de menace entre elles et à quantifier l'évolution temporelle de leur pertinence. Nous avons implémenté deux mécanismes complémentaires qui répondent à ces exigences : une topologie de graphe orienté pour la modélisation multi-sauts des chaînes d'attaque, et un moteur de déclin paramétré pour la gestion algorithmique du cycle de vie des indicateurs de compromission.

\subsection{Topologie Graphe (Neo4j) et Modélisation Multi-Sauts}

\subsubsection{Le modèle de données : Graphe de propriétés labelé}

Comme justifié au Chapitre~II (section~II.3.c), Neo4j est retenu pour modéliser les relations complexes entre entités de renseignement. Le modèle de données sous-jacent est le \textbf{graphe de propriétés labelé} (\textit{Labeled Property Graph}, LPG), une structure formellement définie comme un tuple $G = (N, R, L_N, L_R, P)$ où :

\begin{itemize}
    \item $N$ est l'ensemble des n\oe uds, chaque n\oe ud représentant une entité STIX (acteur de menace, campagne, programme malveillant, vulnérabilité, infrastructure, indicateur).
    \item $R \subseteq N \times N$ est l'ensemble des arêtes orientées, chaque arête représentant une relation STIX (\texttt{attributed-to}, \texttt{uses}, \texttt{exploits}, \texttt{targets}, \texttt{indicates}).
    \item $L_N : N \rightarrow 2^{\Lambda}$ est la fonction d'étiquetage qui associe à chaque n\oe ud un ensemble de labels (par exemple, \texttt{ThreatActor}, \texttt{Campaign}).
    \item $L_R : R \rightarrow \Lambda$ est la fonction d'étiquetage des arêtes (par exemple, \texttt{ATTRIBUTED\_TO}).
    \item $P$ est l'ensemble des propriétés clé-valeur attachées aux n\oe uds et aux arêtes (identifiant STIX, nom, score de confiance, horodatage).
\end{itemize}

Ce modèle permet de représenter nativement les 18 types d'objets STIX~2.1 comme des n\oe uds labelés et les relations STIX comme des arêtes orientées, sans la transformation structurelle (\textit{impedance mismatch}) qu'imposerait un schéma relationnel.

\subsubsection{Contiguïté sans index (\textit{Index-Free Adjacency})}

La propriété architecturale fondamentale qui distingue Neo4j des bases de données relationnelles pour les requêtes de traversée est la \textbf{contiguïté sans index} (\textit{index-free adjacency}). Dans une base relationnelle, la résolution d'une jointure entre deux tables nécessite une recherche par index (typiquement un arbre B+ de complexité $\mathcal{O}(\log n)$ par clé étrangère). Pour une traversée de $k$ jointures sur une table de $n$ enregistrements, la complexité totale est de $\mathcal{O}(k \cdot \log n)$ dans le cas indexé, et de $\mathcal{O}(k \cdot n)$ dans le cas non indexé --- cette complexité croissant donc avec la taille totale de la base.

Dans Neo4j, chaque n\oe ud stocke physiquement des pointeurs directs vers ses n\oe uds adjacents. La traversée d'une arête est une opération de déréférencement de pointeur en $\mathcal{O}(1)$, indépendamment du nombre total de n\oe uds dans le graphe. Pour une traversée de $k$ sauts visitant en moyenne $b$ voisins par n\oe ud (\textit{branching factor}), la complexité totale est de $\mathcal{O}(b^k)$ --- elle dépend de la topologie locale du sous-graphe exploré, mais \textit{pas} de la taille globale du graphe. Cette propriété est déterminante pour une plateforme de renseignement : à mesure que la base de connaissances s'enrichit en millions d'indicateurs, le temps de traversée d'une chaîne d'attaque spécifique reste constant, car il ne dépend que de la profondeur de la chaîne et du degré moyen des n\oe uds impliqués.

\subsubsection{Requêtes de corrélation multi-sauts en Cypher}

Le langage de requête Cypher, natif de Neo4j, exprime les traversées de graphe sous une forme déclarative qui reflète directement la topologie recherchée. Considérons le scénario opérationnel suivant : un analyste identifie une adresse IP suspecte et souhaite reconstruire l'intégralité de la chaîne d'attaque associée --- de l'acteur responsable jusqu'aux infrastructures ciblées, en passant par les campagnes, les programmes malveillants et les vulnérabilités exploitées. Cette reconstruction, que nous désignons par le terme de \textbf{rayon d'impact} (\textit{Blast Radius}) d'un indicateur, s'exprime en une seule requête Cypher :

\begin{lstlisting}[language=, caption={Requête Cypher de corrélation multi-sauts : reconstruction du rayon d'impact d'un indicateur.}, label={lst:cypher}]
// Rayon d'impact (Blast Radius) d'un indicateur
// Traversee : Indicateur -> Malware -> Campagne
//          -> Acteur -> Vulnerabilite -> Infrastructure
MATCH (ioc:Indicator {value: '185.220.101.45'})
MATCH path = (ioc)-[:INDICATES]->(m:Malware)
              -[:USED_BY]->(c:Campaign)
              -[:ATTRIBUTED_TO]->(a:ThreatActor)
OPTIONAL MATCH vuln_path = (m)-[:EXPLOITS]->(v:Vulnerability)
OPTIONAL MATCH infra_path = (c)-[:TARGETS]->(i:Infrastructure)
RETURN
  a.name        AS acteur,
  a.aliases     AS alias_connus,
  c.name        AS campagne,
  c.first_seen  AS debut_campagne,
  m.name        AS malware,
  m.malware_types AS type_malware,
  v.name        AS vulnerabilite,
  v.cvss_score  AS score_cvss,
  i.name        AS infrastructure_ciblee,
  length(path)  AS profondeur_traversee
ORDER BY v.cvss_score DESC
\end{lstlisting}

Cette requête unique reconstruit, en une seule traversée, la chaîne complète : \textbf{Indicateur $\rightarrow$ Programme malveillant $\rightarrow$ Campagne $\rightarrow$ Acteur de menace}, avec des branches optionnelles vers les \textbf{Vulnérabilités exploitées} et les \textbf{Infrastructures ciblées}. Dans une base relationnelle classique, cette reconstruction nécessiterait quatre à cinq jointures successives (\texttt{JOIN}) entre des tables distinctes, dont le coût computationnel croîtrait avec le volume de chaque table. Grâce à la contiguïté sans index de Neo4j, cette opération s'exécute en temps proportionnel au nombre de relations traversées ($k = 4$ à $5$ sauts), indépendamment du fait que la base contienne $10^3$ ou $10^7$ n\oe uds.

\subsubsection{Traversée programmatique bornée via MongoDB \$graphLookup}

Pour les traversées ne nécessitant pas la pleine expressivité de Cypher --- par exemple, la découverte de tous les n\oe uds atteignables depuis un objet STIX donné dans un rayon de $k$ sauts --- nous exploitons l'opérateur d'agrégation \texttt{\$graphLookup} de MongoDB. Cet opérateur effectue une recherche récursive en profondeur dans la collection \texttt{stix\_relationships}, paramétrable jusqu'à $k-1$ niveaux de profondeur via le champ \texttt{maxDepth}. Le résultat est un document enrichi contenant l'ensemble des relations et des n\oe uds découverts, avec un champ \texttt{depth} indiquant la distance de chaque n\oe ud par rapport à la racine.

Cette double stratégie --- Neo4j pour les analyses exploratoires complexes à profondeur variable, MongoDB pour les traversées programmatiques bornées et les opérations de masse --- optimise le rapport entre expressivité et performance en acheminant chaque catégorie de requête vers le moteur de stockage le plus adapté à sa nature.

\subsection{Le Moteur de Déclin (\textit{Decay Engine})}

Un indicateur de compromission ne conserve pas indéfiniment sa pertinence opérationnelle. La dynamique des infrastructures d'attaque impose une réalité asymétrique : une adresse IP utilisée comme serveur de commande et contrôle sur une instance éphémère d'infrastructure en nuage peut être réassignée à un hébergeur légitime en quelques heures, tandis qu'une empreinte SHA-256 d'un binaire malveillant reste pertinente pour la détection forensique pendant plusieurs années. Nous avons formalisé cette réalité par un moteur de déclin paramétré, inspiré des modèles de décroissance de la plateforme MISP (\textit{MISP Decaying Models}) et étendu avec des paramètres contextuels spécifiques à chaque type et sous-type d'artefact.

\subsubsection{Modèle mathématique de décroissance exponentielle}

Nous modélisons la dégradation temporelle de la pertinence d'un indicateur par une fonction de décroissance exponentielle paramétrée par la demi-vie de l'artefact. Pour un indicateur de type $\tau$ observé pour la dernière fois il y a un intervalle temporel $\Delta t$, le score de déclin $D(\Delta t, \tau)$ est défini par :

\begin{equation}
\label{eq:decay}
D(\Delta t, \tau) = \exp\!\left( -\frac{\ln 2 \cdot \Delta t}{t_{1/2}(\tau)} \right)
\end{equation}

où $t_{1/2}(\tau)$ désigne la demi-vie spécifique au type $\tau$, c'est-à-dire l'intervalle temporel au bout duquel le score de déclin atteint exactement $0{,}5$. La constante $\ln 2 \approx 0{,}6931$ est intégrée au numérateur de l'exposant pour garantir cette propriété par construction mathématique. Nous démontrons les trois propriétés fondamentales du modèle :

\begin{itemize}
    \item \textbf{Fraîcheur maximale à l'instant d'observation} : $D(0, \tau) = \exp(0) = 1{,}0$ --- un indicateur fraîchement observé possède une pertinence maximale.
    \item \textbf{Demi-vie exacte} : $D(t_{1/2}, \tau) = \exp(-\ln 2) = 2^{-1} = 0{,}5$ --- la pertinence est exactement réduite de moitié après une demi-vie, quelle que soit la valeur de $t_{1/2}$.
    \item \textbf{Convergence asymptotique} : $\displaystyle\lim_{\Delta t \to \infty} D(\Delta t, \tau) = 0$ --- la pertinence tend asymptotiquement vers zéro sans jamais l'atteindre, conformément au comportement physique attendu.
\end{itemize}

L'implémentation est réduite à une fonction pure, sans état interne, appelable depuis les processus de travail asynchrones comme depuis le code synchrone :

\begin{lstlisting}[language=Python, caption={Fonction de décroissance exponentielle paramétrée par la demi-vie.}, label={lst:decay-func}]
import math

# ln(2) : garantit decay(t=half_life) == 0.5 exactement
_LN2 = 0.6931471805599453

def exponential_decay(
    delta_time: float,
    half_life: float,
) -> float:
    """
    Calcul du score de declin.
    delta_time et half_life doivent etre exprimes
    dans la meme unite temporelle (heures ou jours).
    Retourne un score dans [0.0, 1.0].
    """
    return math.exp(-_LN2 * delta_time / half_life)
\end{lstlisting}

\subsubsection{Demi-vies paramétrées par type d'artefact}

Nous avons calibré les demi-vies de notre moteur à partir de l'analyse opérationnelle des cycles de vie des indicateurs, en croisant les données publiées par les équipes de réponse à incident (CERT-FR, US-CERT) et les observations empiriques de notre propre collecte. Le Tableau~\ref{tab:decay-half-lives} présente les valeurs retenues avec leur justification technique.

\begin{table}[ht]
\centering
\caption{Demi-vies paramétrées par type d'indicateur de compromission.}
\label{tab:decay-half-lives}
\begin{tabular}{|l|l|p{7.5cm}|}
\hline
\textbf{Type d'artefact} & \textbf{Demi-vie $t_{1/2}$} & \textbf{Justification technique} \\
\hline
IP — Cloud légitime (AWS, Azure, GCP) & 6 heures & Les fournisseurs d'infrastructure en nuage recyclent les adresses IP élastiques en quelques heures après libération de l'instance. \\
\hline
IP — Proxy résidentiel & 24 heures & Les pools résidentiels sont renouvelés quotidiennement par les opérateurs de botnets résidentiels. \\
\hline
IP — N\oe ud de sortie Tor & 48 heures & Les relais Tor suivent un calendrier de rotation régulier dicté par les directives du consensus. \\
\hline
IP — Hébergement blindé (\textit{bulletproof}) & 72 heures & L'opérateur contrôle l'infrastructure dédiée; la rotation est délibérément lente pour maintenir la stabilité du C2. \\
\hline
IP — Serveur dédié criminel & 96 heures & Infrastructure louée sous fausse identité; durée de vie prolongée tant que l'hébergeur ne reçoit pas de signalement. \\
\hline
Domaine — TLD à fort abus (.tk, .ml, .ga) & 3--4 jours & TLD gratuits fréquemment réenregistrés en masse via des services automatisés d'enregistrement. \\
\hline
Domaine — TLD standard (.com, .net) & 25--30 jours & Rotation modérée; le domaine peut être saisi judiciairement ou expirer selon le registraire. \\
\hline
URL (C2, hameçonnage) & 48 heures & Les chemins de rappel C2 et les pages de hameçonnage sont modifiés à haute fréquence pour échapper aux listes de blocage. \\
\hline
SHA-256 & 730 jours & Empreinte cryptographiquement liée au binaire; pertinence forensique durable. Plancher à $0{,}30$. \\
\hline
SHA-1 & 548 jours & Empreinte moins robuste (vulnérabilité aux collisions) mais encore exploitable en détection. \\
\hline
MD5 & 365 jours & Risque de collision démontré (attaque de Wang et al., 2004); poids opérationnel réduit. \\
\hline
Empreinte JA3/JA4 & 180 jours & Liée à l'implémentation TLS du programme malveillant; persiste jusqu'à la mise à jour majeure du code. \\
\hline
CVE & 180--365 jours & Décroît en probabilité d'exploitation active une fois le taux de correction supérieur à 80~\%. Plancher à $0{,}10$. \\
\hline
Courriel (expéditeur d'hameçonnage) & 14 jours & Les adresses d'envoi sont jetables et changent à chaque vague de campagne. \\
\hline
Courriel (enregistrement C2) & 90 jours & L'acteur conserve l'adresse pour le renouvellement de l'infrastructure (certificats, domaines). \\
\hline
\end{tabular}
\end{table}

Nous soulignons deux raffinements importants dans notre implémentation. Premièrement, les empreintes cryptographiques (SHA-256, SHA-1, MD5) possèdent un \textbf{plancher de déclin} (\textit{decay floor}) fixé à $0{,}30$ : même après plusieurs années, une empreinte de fichier malveillant ne descend jamais en dessous de ce seuil, car elle conserve sa valeur pour la détection forensique rétrospective. Deuxièmement, les demi-vies des domaines sont \textbf{modulées par le tier du registraire} via un facteur multiplicatif : un registraire blindé (\textit{bulletproof}, tier~1) double la demi-vie (facteur $2{,}0$), car l'acteur contrôle le domaine et ne le laissera pas expirer; un registraire réputé (tier~4, par exemple GoDaddy ou Namecheap) divise la demi-vie par deux (facteur $0{,}5$), car le domaine peut être suspendu ou transféré rapidement suite à un signalement d'abus.

\subsubsection{Discrétisation en états opérationnels}

Le score continu $D \in [0, 1]$ est informatiquement précis mais opérationnellement inadapté à la prise de décision automatisée dans les SIEM et les plateformes de réponse à incident. Nous avons donc défini une fonction de discrétisation qui projette le score continu vers quatre états opérationnels, chacun associé à une action défensive prescrite :

\begin{itemize}
    \item \textbf{VALID} ($D > 0{,}80$) : l'indicateur est considéré comme actif et pertinent. \textit{Action prescrite} : blocage immédiat dans les règles de pare-feu et de proxy, génération d'alerte prioritaire dans le SIEM.
    \item \textbf{DEGRADING} ($0{,}50 < D \leq 0{,}80$) : l'indicateur est en cours de dégradation. \textit{Action prescrite} : maintien du blocage avec demande de validation manuelle par un analyste avant renouvellement.
    \item \textbf{STALE} ($0{,}15 < D \leq 0{,}50$) : l'indicateur est périmé. \textit{Action prescrite} : levée du blocage actif, conservation en liste de surveillance passive. Vérification recommandée auprès d'une source externe avant tout réengagement.
    \item \textbf{OBSOLETE} ($D \leq 0{,}15$) : l'indicateur est considéré comme obsolète. \textit{Action prescrite} : retrait complet des règles de blocage. Archivage en base historique pour analyse rétrospective uniquement.
\end{itemize}

Deux états supplémentaires, \textbf{CONTESTED} et \textbf{RETRACTED}, sont attribués manuellement par l'analyste lorsque l'attribution d'un indicateur est contestée par une source fiable ou que l'émetteur original a rétracté l'indicateur. Ces états court-circuitent le score temporel : un indicateur rétracté reçoit un score forcé de $0{,}05$ quelle que soit sa fraîcheur, empêchant ainsi tout blocage fondé sur un renseignement invalidé.

Les seuils de discrétisation ($0{,}80$, $0{,}50$, $0{,}15$) ont été calibrés pour minimiser simultanément deux risques opérationnels antagonistes : le risque de \textit{sous-blocage} (un indicateur encore actif est prématurément retiré des règles défensives) et le risque de \textit{sur-blocage} (un indicateur obsolète maintenu dans les règles produit des faux positifs qui saturent la capacité d'analyse du SOC). L'automatisation de ce cycle de vie --- de l'insertion avec un score de $1{,}0$ jusqu'au retrait automatique sous le seuil de $0{,}15$ --- réduit la charge opérationnelle des analystes SOC, qui n'interviennent manuellement que dans la zone intermédiaire DEGRADING où le jugement humain reste nécessaire.

\subsubsection{Recalcul périodique et boucle d'apprentissage}

Le recalcul des scores de déclin est orchestré par un processus de travail asynchrone dédié, déclenché toutes les heures par l'ordonnanceur Celery Beat (Chapitre~II, section~II.2.g). Ce processus itère sur l'ensemble des indicateurs actifs, recalcule leur score $D$ en fonction du temps écoulé depuis leur dernière observation, met à jour l'état opérationnel et émet un événement de changement d'état vers le bus Redis pour notification en temps réel aux tableaux de bord.

En parallèle, une boucle d'apprentissage (\textit{decay learning loop}) observe le comportement réel des indicateurs --- réapparition dans de nouvelles sources, confirmation par des observations multiples, rétractation --- et ajuste les demi-vies par acteur. Si un groupe adverse spécifique utilise systématiquement ses adresses IP pendant 12 heures avant rotation (au lieu des 36 heures par défaut), le système apprend ce profil et applique une demi-vie personnalisée de 12 heures aux futurs indicateurs attribués à cet acteur. Ce mécanisme de rétroaction transforme le moteur de déclin d'un modèle statique en un système adaptatif qui affine ses paramètres au fil de l'accumulation de données opérationnelles.



\section{Tolérance aux pannes et Orchestration Asynchrone}

La fiabilité d'une plateforme de renseignement opérationnel se mesure à sa capacité à maintenir la disponibilité du service sous charge soutenue --- typiquement $10^3$ indicateurs de compromission ingérés par seconde lors d'une vague de publication simultanée par les flux OSINT --- tout en absorbant les défaillances individuelles de composants sans perte de données. Nous avons conçu l'architecture d'ONYX selon un paradigme de tolérance aux pannes (\textit{Fail-Safe}) fondé sur la ségrégation stricte entre les opérations liées aux entrées-sorties (\textit{I/O-bound}) et les opérations liées au calcul intensif (\textit{CPU-bound}). Nous démontrons dans cette section que cette ségrégation garantit l'absence de point unique de défaillance sur le chemin critique du traitement.

\subsection{Modèle de concurrence I/O-Bound : FastAPI, Uvicorn et le protocole ASGI}

\subsubsection{Le protocole ASGI et la boucle événementielle}

Comme justifié au Chapitre~II (section~II.2.a), le serveur d'application repose sur FastAPI exécuté par Uvicorn. Uvicorn implémente le protocole ASGI (\textit{Asynchronous Server Gateway Interface}), successeur asynchrone du protocole WSGI historique. Là où WSGI traite chaque requête de manière synchrone et séquentielle --- un fil d'exécution par requête, bloqué pendant toute la durée du traitement --- ASGI expose une interface fondée sur des coroutines Python (\texttt{async def}) qui peuvent être suspendues et reprises au gré des opérations d'entrée-sortie.

Le moteur d'exécution sous-jacent est la boucle événementielle \texttt{uvloop} (implémentation C de \texttt{asyncio}), qui pilote un multiplexeur d'entrées-sorties via l'appel système \texttt{epoll} (Linux) ou \texttt{kqueue} (macOS/BSD). Ce multiplexeur surveille simultanément l'ensemble des descripteurs de fichier ouverts --- sockets réseau vers les clients HTTP, connexions vers MongoDB, Elasticsearch, Redis, et sockets WebSocket --- et notifie la boucle événementielle dès qu'un descripteur est prêt pour une opération de lecture ou d'écriture. Ce mécanisme permet à un processus unique de gérer simultanément des milliers de connexions concurrentes sans recourir au multi-threading, éliminant ainsi les problèmes classiques de verrous (\textit{locks}), de conditions de concurrence (\textit{race conditions}) et d'interblocages (\textit{deadlocks}).

\subsubsection{Architecture à blocage nul (\textit{Zero-Blocking Architecture})}

Le principe fondamental de notre architecture I/O-bound est le suivant : \textbf{aucune opération d'entrée-sortie réseau ne bloque le fil d'exécution principal}. Chaque interaction avec un service externe est implémentée comme une coroutine asynchrone qui cède le contrôle à la boucle événementielle via le mot-clé \texttt{await} pendant l'attente de la réponse.

Considérons le flux de traitement d'une requête API typique --- la recherche d'un indicateur de compromission par valeur :

\begin{enumerate}
    \item Uvicorn reçoit la requête HTTP et l'achemine vers le gestionnaire FastAPI correspondant (\texttt{async def get\_ioc(value: str)}).
    \item Le gestionnaire émet une requête vers Elasticsearch (\texttt{await es.search(...)}) pour la recherche textuelle. Pendant l'attente de la réponse réseau (typiquement 2 à 15~ms), la coroutine est suspendue et la boucle événementielle peut traiter d'autres requêtes entrantes.
    \item À la réception de la réponse Elasticsearch, la coroutine reprend, enrichit le résultat par une interrogation de Redis (\texttt{await redis.cache\_get(...)}) pour récupérer le score de déclin en cache. Nouvelle suspension pendant l'aller-retour réseau ($<1$~ms pour Redis in-memory).
    \item Le résultat complet est sérialisé en JSON via ORJSON (sérialiseur compilé en Rust, 3 à 10 fois plus rapide que le module \texttt{json} standard) et retourné au client.
\end{enumerate}

Durant ce flux, le fil d'exécution principal n'est jamais bloqué par une attente réseau. Le temps total de blocage effectif est limité aux opérations CPU pures --- sérialisation JSON, validation Pydantic, construction de la réponse --- qui représentent typiquement moins de 5~\% du temps de traitement total d'une requête.

Notre implémentation du cycle de vie applicatif (\texttt{lifespan}) illustre cette architecture à plus grande échelle : au démarrage, les services de base de données sont initialisés de manière asynchrone (\texttt{await es.initialize()}, \texttt{await mongo.initialize()}, \texttt{await redis.initialize()}), puis cinq tâches d'ingestion en temps réel sont lancées comme coroutines concurrentes via \texttt{asyncio.create\_task()} --- le poller OSINT, l'ingesteur géopolitique, le moteur de rapports dynamiques NLP, la boucle de recalcul des scores de déclin et la boucle d'apprentissage des demi-vies. Chacune de ces tâches opère de manière indépendante dans la boucle événementielle, sans bloquer le traitement des requêtes API entrantes, et sans contention de verrous entre elles.

\subsection{Délégation CPU-Bound : Celery Workers et Redis comme Courtier de Messages}

Les opérations de calcul intensif --- inférence du modèle SciBERT (vectorisation d'un rapport en $\mathbb{R}^{768}$), traversée de graphes profonds dans Neo4j, recalcul de masse des scores de déclin sur $10^5$ indicateurs --- ne peuvent pas être exécutées dans la boucle événementielle sans conséquences critiques. Une opération CPU-bound de 500~ms bloquerait le fil unique d'Uvicorn, suspendant le traitement de toutes les requêtes API concurrentes et des WebSockets actifs pendant cette durée. À un débit de 200 requêtes par seconde, cela produirait une file d'attente de 100 requêtes en attente --- un effondrement de performance inacceptable pour un usage opérationnel.

Nous avons résolu cette contrainte par la délégation systématique des charges CPU-bound vers des processus de travail Celery (\textit{Celery Workers}) exécutés dans des processus distincts, avec Redis comme courtier de messages (\textit{Message Broker}).

\subsubsection{Redis comme courtier de messages in-memory}

Redis opère intégralement en mémoire volatile avec une politique d'éviction LRU (\textit{Least Recently Used}) et sert simultanément trois fonctions architecturales dans notre plateforme :

\begin{enumerate}
    \item \textbf{Courtier de messages Celery} : Redis reçoit les tâches soumises par l'API sous forme de messages sérialisés en JSON, les stocke dans quatre files d'attente nommées et ségrégées par domaine fonctionnel (\texttt{crawlers}, \texttt{nlp}, \texttt{analysis}, \texttt{default}), et les distribue aux processus de travail abonnés à chaque file selon une politique FIFO (\textit{First In, First Out}). La latence d'enfilement est inférieure à 0{,}1~ms grâce au stockage in-memory, ce qui garantit que la soumission d'une tâche depuis l'API ne produit aucun blocage perceptible.
    \item \textbf{Stockage des résultats} : les résultats des tâches asynchrones sont persistés dans Redis avec un temps d'expiration configurable, permettant au client de récupérer le résultat à tout moment après la soumission via l'identifiant de tâche Celery.
    \item \textbf{Bus d'événements temps réel} : les Redis Streams servent de bus de publication-abonnement (\textit{pub/sub}) pour la diffusion des événements système (nouvelles détections d'IoC, résultats d'extraction NLP, alertes géopolitiques, changements d'état de déclin) vers les consommateurs WebSocket connectés aux tableaux de bord.
\end{enumerate}

Le routage des tâches est implémenté par espace de noms : toutes les tâches du module \texttt{onyx\_crawlers} sont automatiquement acheminées vers la file \texttt{crawlers}, les tâches \texttt{onyx\_nlp} vers la file \texttt{nlp} et les tâches \texttt{onyx\_analyzers} vers la file \texttt{analysis}. Cette ségrégation par file garantit qu'une surcharge de tâches de collecte --- par exemple, lors du lancement simultané de 50 crawlers Dark Web --- ne bloque jamais les tâches d'analyse NLP, et inversement. Chaque file peut être dimensionnée indépendamment par ajout de processus de travail dédiés, assurant l'extensibilité horizontale exigée par l'exigence non fonctionnelle ENF-3 du Chapitre~I.

L'ordonnanceur périodique Celery Beat complète cette architecture en déclenchant automatiquement les tâches récurrentes : surveillance de Pastebin toutes les 5 minutes, surveillance de GitHub toutes les 10 minutes, recalcul des scores de déclin toutes les heures et rafraîchissement du cache du tableau de bord toutes les 30 secondes.

\subsubsection{Tolérance aux pannes par acquittement tardif (\texttt{task\_acks\_late})}

La configuration \texttt{task\_acks\_late = True} de Celery constitue la pierre angulaire de notre stratégie de tolérance aux pannes pour les tâches distribuées. Dans le mode par défaut (\texttt{task\_acks\_late = False}), un message est acquitté dès sa réception par le processus de travail, \textit{avant} l'exécution de la tâche. Si le processus subit une terminaison imprévue --- arrêt brutal du conteneur Docker, erreur de segmentation dans une bibliothèque C native, épuisement mémoire par le modèle SciBERT --- le message est définitivement perdu, car le courtier le considère comme traité.

Avec \texttt{task\_acks\_late = True}, le message n'est acquitté qu'après l'exécution complète et réussie de la tâche. En cas de terminaison imprévue du processus de travail, le message non acquitté est automatiquement réinjecté dans la file d'attente par Redis et redistribué à un autre processus de travail disponible. Ce mécanisme garantit la propriété \textit{at-least-once delivery} : chaque tâche est exécutée au moins une fois, même en cas de défaillance matérielle d'un n\oe ud de calcul.

Le paramètre complémentaire \texttt{worker\_prefetch\_multiplier = 1} empêche chaque processus de travail de pré-extraire plus d'un message à la fois, limitant ainsi le nombre de tâches potentiellement perdues en cas de crash à exactement une --- qui sera de toute façon redistribuée grâce à l'acquittement tardif.

\subsubsection{Résilience réseau par tentatives exponentielles (\textit{Exponential Backoff})}

Les interruptions transitoires de connectivité réseau entre les services --- redémarrage d'un conteneur MongoDB, reconfiguration réseau d'Elasticsearch, saturation temporaire d'un pool de connexions --- constituent un scénario opérationnel courant dans un environnement conteneurisé. Une stratégie de reconnexion naïve (tentative immédiate, échec, tentative immédiate, échec) produit un phénomène d'avalanche (\textit{thundering herd}) : tous les clients tentent de se reconnecter simultanément, saturant le service à peine rétabli.

Nous avons implémenté une stratégie de tentatives exponentielles (\textit{exponential backoff}) via la bibliothèque \texttt{tenacity}, appliquée à toutes les initialisations de services critiques. La configuration retenue est la suivante :

\begin{lstlisting}[language=Python, caption={Résilience réseau par tentatives exponentielles via \texttt{tenacity}.}, label={lst:backoff}]
from tenacity import (
    retry,
    stop_after_attempt,
    wait_exponential,
)

@retry(
    stop=stop_after_attempt(5),
    wait=wait_exponential(multiplier=2, max=30),
)
async def initialize(self) -> None:
    """
    Initialisation du service MongoDB avec resilience.
    Strategie : 5 tentatives, delai exponentiel.

    Sequence de delais :
      Tentative 1 : echec -> attente 2s
      Tentative 2 : echec -> attente 4s
      Tentative 3 : echec -> attente 8s
      Tentative 4 : echec -> attente 16s
      Tentative 5 : echec -> attente 30s (plafond)

    Duree totale maximale avant abandon : 60s.
    """
    info = await self._client.server_info()
    logger.info(
        "MongoDB connected: version=%s",
        info.get("version"),
    )
\end{lstlisting}

Le délai entre la $k$-ième tentative et la suivante est calculé par $d_k = \min(2^k, 30)$ secondes, produisant la séquence $\{2, 4, 8, 16, 30\}$. La durée totale maximale avant abandon est de 60 secondes, ce qui absorbe la très grande majorité des interruptions transitoires (redémarrage d'un conteneur : 5 à 15 secondes; reconfiguration réseau : 10 à 30 secondes) tout en garantissant une détection rapide des défaillances permanentes.

\subsubsection{Dégradation gracieuse et isolation des défaillances}

Notre architecture garantit la propriété de \textbf{dégradation gracieuse} (\textit{graceful degradation}) : en cas d'indisponibilité d'un service non critique, la plateforme continue de fonctionner en mode dégradé plutôt que de s'arrêter entièrement. L'initialisation de chaque service au démarrage est encapsulée dans un bloc \texttt{try/except} indépendant : si Qdrant est indisponible, les fonctionnalités de recherche sémantique sont désactivées mais l'API REST, la recherche textuelle Elasticsearch et l'ingestion d'indicateurs continuent de fonctionner. Si le moteur de géolocalisation GeoIP ne s'initialise pas, les indicateurs sont ingérés sans enrichissement géographique.

Cette isolation des défaillances est renforcée par la ségrégation des files de tâches Celery : un incident dans le module de collecte (\textit{crawlers}) --- par exemple, l'indisponibilité du proxy Tor --- n'affecte ni l'API, ni le moteur d'analyse NLP, ni les files d'export TAXII. Chaque domaine fonctionnel opère dans un compartiment étanche, connecté aux autres exclusivement par le bus d'événements Redis.

\bigskip
\noindent\rule{\textwidth}{0.4pt}
\bigskip

La conjonction des cinq piliers techniques exposés dans ce chapitre --- acquisition furtive avec sécurité opérationnelle intégrée (section~3.1), pipeline ETL normalisé selon le standard STIX~2.1 (section~3.2), intelligence artificielle sémantique par extraction NER et vectorisation SciBERT (section~3.3), modélisation par graphe de menace avec moteur de déclin paramétré (section~3.4) et orchestration asynchrone tolérante aux pannes (section~3.5) --- forme une chaîne de traitement cohérente qui transforme le signal brut, issu de sources hétérogènes et potentiellement hostiles, en renseignement contextualisé, corrélé et immédiatement exploitable par les équipes de défense. Nous détaillons, dans le chapitre suivant, les interfaces de restitution et les mécanismes de visualisation qui permettent aux analystes d'exploiter ce renseignement : le tableau de bord React, la cartographie géospatiale, l'assistant conversationnel fondé sur l'intelligence artificielle générative et les mécanismes d'exportation vers les écosystèmes SIEM partenaires.
